% GENERATED FILE. Edit canonical lessons, then run npm run generate:books.
% canonical-source-hash: fnv1a64:ffe6e433fa51cef4
% canonical-lessons: SA-C161-sankhah, SA-C161-venuh, SA-C161-mrdangah, SA-C161-nupuram, SA-C161-kundalam

\chapter{Conch, Flute, Mridanga drum, Anklet, Earring}
\label{ch:sa-conch-flute-mridanga-drum-anklet-earring}
\hlchaptermodality{161}

\begin{chapteropening}
\textbf{By the end of this chapter:} \emph{I can say a conch, a flute, a mridanga drum, an anklet and an earring.}

Complete the last lesson of chapter 161: I can say a conch, a flute, a mridanga drum, an anklet and an earring.
\end{chapteropening}

\section[\texorpdfstring{\sk{शंखः} (śaṅkhaḥ)}{शंखः (śaṅkhaḥ)}]{\sk{शंखः} (śaṅkhaḥ) — a conch}

\label{lesson:SA-C161-sankhah}



\begin{quote}
Before the new one: say the Sanskrit for eating, then the Sanskrit for a small pitcher.
\end{quote}

\subsection*{You'll want to know: \sk{शंखः}}
\textbf{\sk{शंखः}} — \emph{śaṅkhaḥ} — \textquotedblleft{}a conch\textquotedblright{}.

\begin{grammarlens}[title={What you've built}]
One more word for this chapter.
\end{grammarlens}

\subsection*{Your turn}
\begin{itemize}
\raggedright
  \item \emph{Say it:} \emph{śaṅkhaḥ}
  \item \emph{Say it:} \emph{śaṅkhaḥ}, once more
  \item \emph{Say it:} say \emph{śaṅkhaḥ}
  \item \emph{Recall:} say the Sanskrit for eating, then the Sanskrit for a small pitcher, then say \emph{śaṅkhaḥ} again
\end{itemize}

\begin{tcolorbox}[breakable,colback=teal!4,colframe=teal!35!black,title={Before you move on}]
What does \sk{शंखः} mean? (\textquotedblleft{}A conch\textquotedblright{}.) Say it once more.
\end{tcolorbox}
\section[\texorpdfstring{\sk{वेणुः} (veṇuḥ)}{वेणुः (veṇuḥ)}]{\sk{वेणुः} (veṇuḥ) — a flute}

\label{lesson:SA-C161-venuh}



\begin{quote}
Before the new one: say the Sanskrit for a small pitcher, then the Sanskrit for a conch.
\end{quote}

\subsection*{You'll want to know: \sk{वेणुः}}
\textbf{\sk{वेणुः}} — \emph{veṇuḥ} — \textquotedblleft{}a flute\textquotedblright{}.

\begin{grammarlens}[title={What you've built}]
One more word for this chapter.
\end{grammarlens}

\subsection*{Your turn}
\begin{itemize}
\raggedright
  \item \emph{Say it:} \emph{veṇuḥ}
  \item \emph{Say it:} \emph{veṇuḥ}, once more
  \item \emph{Say it:} say \emph{veṇuḥ}
  \item \emph{Recall:} say the Sanskrit for a small pitcher, then the Sanskrit for a conch, then say \emph{veṇuḥ} again
\end{itemize}

\begin{tcolorbox}[breakable,colback=teal!4,colframe=teal!35!black,title={Before you move on}]
What does \sk{वेणुः} mean? (\textquotedblleft{}A flute\textquotedblright{}.) Say it once more.
\end{tcolorbox}
\section[\texorpdfstring{\sk{मृदंगः} (mṛdaṅgaḥ)}{मृदंगः (mṛdaṅgaḥ)}]{\sk{मृदंगः} (mṛdaṅgaḥ) — a mridanga drum}

\label{lesson:SA-C161-mrdangah}



\begin{quote}
Before the new one: say the Sanskrit for a conch, then the Sanskrit for a flute.
\end{quote}

\subsection*{You'll want to know: \sk{मृदंगः}}
\textbf{\sk{मृदंगः}} — \emph{mṛdaṅgaḥ} — \textquotedblleft{}a mridanga drum\textquotedblright{}.

\begin{grammarlens}[title={What you've built}]
One more word for this chapter.
\end{grammarlens}

\subsection*{Your turn}
\begin{itemize}
\raggedright
  \item \emph{Say it:} \emph{mṛdaṅgaḥ}
  \item \emph{Say it:} \emph{mṛdaṅgaḥ}, once more
  \item \emph{Say it:} say \emph{mṛdaṅgaḥ}
  \item \emph{Recall:} say the Sanskrit for a conch, then the Sanskrit for a flute, then say \emph{mṛdaṅgaḥ} again
\end{itemize}

\begin{tcolorbox}[breakable,colback=teal!4,colframe=teal!35!black,title={Before you move on}]
What does \sk{मृदंगः} mean? (\textquotedblleft{}A mridanga drum\textquotedblright{}.) Say it once more.
\end{tcolorbox}
\section[\texorpdfstring{\sk{नूपुरम्} (nūpuram)}{नूपुरम् (nūpuram)}]{\sk{नूपुरम्} (nūpuram) — an anklet}

\label{lesson:SA-C161-nupuram}



\begin{quote}
Before the new one: say the Sanskrit for a flute, then the Sanskrit for a mridanga drum.
\end{quote}

\subsection*{You'll want to know: \sk{नूपुरम्}}
\textbf{\sk{नूपुरम्}} — \emph{nūpuram} — \textquotedblleft{}an anklet\textquotedblright{}.

\begin{grammarlens}[title={What you've built}]
One more word for this chapter.
\end{grammarlens}

\subsection*{Your turn}
\begin{itemize}
\raggedright
  \item \emph{Say it:} \emph{nūpuram}
  \item \emph{Say it:} \emph{nūpuram}, once more
  \item \emph{Say it:} say \emph{nūpuram}
  \item \emph{Recall:} say the Sanskrit for a flute, then the Sanskrit for a mridanga drum, then say \emph{nūpuram} again
\end{itemize}

\begin{tcolorbox}[breakable,colback=teal!4,colframe=teal!35!black,title={Before you move on}]
What does \sk{नूपुरम्} mean? (\textquotedblleft{}An anklet\textquotedblright{}.) Say it once more.
\end{tcolorbox}
\section[\texorpdfstring{\sk{कुण्डलम्} (kuṇḍalam)}{कुण्डलम् (kuṇḍalam)}]{\sk{कुण्डलम्} (kuṇḍalam) — an earring}

\label{lesson:SA-C161-kundalam}



\begin{quote}
Before the new one: say the Sanskrit for a mridanga drum, then the Sanskrit for an anklet.
\end{quote}

\subsection*{You'll want to know: \sk{कुण्डलम्}}
\textbf{\sk{कुण्डलम्}} — \emph{kuṇḍalam} — \textquotedblleft{}an earring\textquotedblright{}.

\begin{grammarlens}[title={What you've built}]
One more word for this chapter.
\end{grammarlens}

\subsection*{Your turn}
\begin{itemize}
\raggedright
  \item \emph{Say it:} \emph{kuṇḍalam}
  \item \emph{Say it:} \emph{kuṇḍalam}, once more
  \item \emph{Say it:} say \emph{kuṇḍalam}
  \item \emph{Recall:} say the Sanskrit for a mridanga drum, then the Sanskrit for an anklet, then say \emph{kuṇḍalam} again
\end{itemize}

\begin{tcolorbox}[breakable,colback=teal!4,colframe=teal!35!black,title={Before you move on}]
What does \sk{कुण्डलम्} mean? (\textquotedblleft{}An earring\textquotedblright{}.) Say it once more.
\end{tcolorbox}
